\subsection{代数的类}

我们用 \(\mathrm{Iso}_{\mathrm{K}}\{\mathrm{A},\mathrm{B}\}\) 表示关系

“A 与 B 是同构的有限次数 K-代数。”

这是关于 A 与 B 的一个等价关系。设 A 是一个有限次数的 K-代数；A 的类，记作 \(\operatorname{cl}(A)\)（有时也记作 \(\operatorname{cl}_{\mathrm{K}}(\mathrm{A})\)），是对关系 \(\operatorname{Iso}_{\mathrm{K}}\) 而言与 A 等价的对象的类（集合论, II, §6, No.~9, 第 122 页）。按定义，\(\operatorname{cl}(A)\) 是一个同构于 A 的 K-代数；两个有限次数的 K-代数属于同一类当且仅当它们同构。

我们用 \(\mathcal{A}\) 表示偶 \((\mathrm{W},\mu)\) 所成之集，其中 \(\mathrm{W}\) 是 \(\mathrm{K}^{(\mathbf{N})}\) 的一个在 \(\mathrm{K}\) 上有限维的线性子空间，而 \(\mu\) 是一个从 \(\mathrm{W} \times \mathrm{W}\) 到 \(\mathrm{W}\) 的 K-双线性映射，它使 \(\mathrm{W}\) 成为一个（结合且含幺的）K-代数。每个有限次数的 K-代数都同构于这样一个代数。由出处同前，关系

\[
\alpha \text {   is   a   class   of   K - algebras   of   finite   degree }
\]

因而在 \(\alpha\) 中是集合化的（集合论, II, §1, No.~4, 第 68 页）。我们用 \(\mathcal{C}_{\mathrm{K}}\) 表示有限次数 K-代数的类所成之集。

命题 1. —— 集合 \(C_{K}\)，赋以由 \((\alpha, \beta) \mapsto \operatorname{cl}(\alpha \otimes_{\mathrm{K}} \beta)\) 给出的合成律，是一个交换幺半群。\(C_{K}\) 的单位元素是 K-代数 K 的类 \(\varepsilon\)。此外，若 A 与 B 是有限次数的 K-代数，则我们有关系

\[
\operatorname{cl} (\mathrm{A} \otimes_ {\mathrm{K}} \mathrm{B}) = \operatorname{cl} (\mathrm{A}) \operatorname{cl} (\mathrm{B}).\tag{1}
\]

设 A、B、C 是有限次数的 K-代数，其类分别为 \(\alpha\)、\(\beta\) 与 \(\gamma\)。K-代数 A 与 \(\alpha\) 同构，B 与 \(\beta\) 亦然。因此，K-代数 \(A \otimes_{K} B\) 与 \(\alpha \otimes_{K} \beta\) 同构，且我们有 \(\operatorname{cl}(A \otimes_{K} B) = \operatorname{cl}(\alpha \otimes_{K} \beta)\)，这就给出公式 (1)。由此得出，\((\alpha\beta)\gamma\) 是 K-代数 \((A \otimes_{K} B) \otimes_{K} C\) 的类，而 \(\alpha(\beta\gamma)\) 是 K-代数 \(A \otimes_{K} (B \otimes_{K} C)\) 的类。而这两个 K-代数是同构的（III, §4, No.~1, 第 461 页），故我们有等式 \((\alpha\beta)\gamma = \alpha(\beta\gamma)\)。类似地，关系 \(\alpha\varepsilon = \varepsilon\alpha = \alpha\) 由 K-代数 \(A \otimes_{K} K\)、\(K \otimes_{K} A\) 与 A 同构这一事实得出，而关系 \(\alpha\beta = \beta\alpha\) 由代数 \(A \otimes_{K} B\) 与 \(B \otimes_{K} A\) 同构这一事实得出。

在集合 \(C_{K}\) 中，关系“\(\alpha\) 与 \(\beta\) 是森田等价的代数”是一个等价关系（VIII, 第 100 页）。由 VIII, 第 111 页的命题 13, d)，它与 \(C_{K}\) 上的合成律相容。我们用 \(M_{K}\) 表示 \(C_{K}\) 模这一等价关系的商幺半群，用 \(\varphi\) 表示从 \(C_{K}\) 到 \(M_{K}\) 的典范同态。对每个有限次数的 K-代数 A，我们用 {[}A{]} 表示 \(\operatorname{cl}(A)\) 在 \(\varphi\) 下的像。若 A 与 B 是有限次数的 K-代数，则我们有 {[}A{]} = {[}B{]} 当且仅当 K-代数 A 与 B 是森田等价的；此外，我们有关系

\[
[ \mathrm{A} \otimes_ {\mathrm{K}} \mathrm{B} ] = [ \mathrm{A} ] [ \mathrm{B} ]\tag{2}
\]

在幺半群 \(M_{K}\) 中。

引理 1. —— 设 A 是一个有限次数的 K-代数，D 是 K 上一个有限次数的体。我们有 \([A]=[D]\)（在 \(M_{K}\) 中）当且仅当存在一个整数 \(n\geqslant1\)，使得 A 同构于 \(\mathbf{M}_{n}(\mathrm{K})\)。

按定义，我们有 \([A]=[D]\) 当且仅当存在一个可逆的 \((A,D)\)-双模，即（VIII, 第 101 页，定理 1）D 上一个非零有限维的右向量空间 V，带有一个 K-代数同构 \(A\to\operatorname{End}_{\mathrm{D}}(\mathrm{V})\)。由此即得引理 1。

